\documentclass[conference]{IEEEtran}

% ─── Packages ─────────────────────────────────────────────────────────────────
\usepackage[utf8]{inputenc}
\usepackage[T1]{fontenc}
\usepackage{amsmath}
\usepackage{amssymb}
\usepackage{graphicx}
\usepackage{cite}
\usepackage{hyperref}
\usepackage{listings}
\usepackage{xcolor}
\usepackage{booktabs}
\usepackage{algorithm}
\usepackage{algorithmic}
\usepackage{tikz}
\usepackage{pgfplots}
\pgfplotsset{compat=1.18}

% ─── Code Style ───────────────────────────────────────────────────────────────
\lstset{
  basicstyle=\ttfamily\footnotesize,
  keywordstyle=\color{blue}\bfseries,
  commentstyle=\color{gray}\itshape,
  stringstyle=\color{orange},
  backgroundcolor=\color{gray!8},
  frame=single,
  breaklines=true,
  numbers=left,
  numberstyle=\tiny\color{gray},
  tabsize=2,
}

% ─── Title & Authors ──────────────────────────────────────────────────────────
\title{Deep Learning-Based Anomaly Detection in\\
       Network Traffic: A Comparative Study}

\author{
  \IEEEauthorblockN{Ahmad Fauzi\textsuperscript{1},
                    Siti Rahayu\textsuperscript{2},
                    Budi Santoso\textsuperscript{3}}
  \IEEEauthorblockA{
    \textsuperscript{1}\textit{Dept. of Computer Science},
    \textit{Universitas Indonesia}, Jakarta, Indonesia\\
    \textsuperscript{2}\textit{Dept. of Electrical Engineering},
    \textit{Institut Teknologi Bandung}, Bandung, Indonesia\\
    \textsuperscript{3}\textit{Dept. of Informatics},
    \textit{Universitas Gadjah Mada}, Yogyakarta, Indonesia\\
    \textsuperscript{1}ahmad.fauzi@ui.ac.id \quad
    \textsuperscript{2}s.rahayu@itb.ac.id \quad
    \textsuperscript{3}b.santoso@ugm.ac.id
  }
}

% ══════════════════════════════════════════════════════════════════════════════
\begin{document}
% ══════════════════════════════════════════════════════════════════════════════

\maketitle

% ─────────────────────────────────────────────────────────────────────────────
\begin{abstract}
	% ─────────────────────────────────────────────────────────────────────────────
	Network intrusion detection remains a critical challenge in modern
	cybersecurity. Traditional signature-based systems fail to detect
	novel or zero-day attacks, motivating the use of machine learning
	approaches. In this paper, we present a comparative study of three
	deep learning architectures --- Long Short-Term Memory (LSTM),
	Convolutional Neural Network (CNN), and Transformer-based models ---
	applied to the problem of anomaly detection in network traffic. We
	evaluate all models on the CICIDS-2017 and NSL-KDD benchmark datasets.
	Our proposed Transformer-based model achieves a detection accuracy of
	\textbf{98.7\%} with a false positive rate of \textbf{0.8\%},
	outperforming existing state-of-the-art methods. We further analyze
	computational trade-offs relevant to deployment in resource-constrained
	edge environments.
\end{abstract}

\begin{IEEEkeywords}
	anomaly detection, deep learning, intrusion detection system,
	LSTM, Transformer, network security, edge computing
\end{IEEEkeywords}

% ─────────────────────────────────────────────────────────────────────────────
\section{Introduction}
\label{sec:intro}
% ─────────────────────────────────────────────────────────────────────────────

The rapid growth of internet-connected devices and services has
dramatically expanded the attack surface available to malicious
actors~\cite{ref:cisco2023}. Intrusion Detection Systems (IDS) serve
as a key defensive layer, monitoring network flows for suspicious
behaviour. Classical IDS rely on hand-crafted signatures and rules
that require continuous manual maintenance and are inherently blind
to previously unseen attack vectors~\cite{ref:liao2013}.

Machine learning has emerged as a compelling alternative, enabling
systems to learn the statistical fingerprint of benign traffic and
flag deviations without requiring explicit rule authoring. Among
modern architectures, deep neural networks have delivered the strongest
empirical results on public benchmarks~\cite{ref:ring2019}.

\subsection{Motivation}

Despite impressive results in controlled evaluations, deploying deep
learning IDS in practice raises several open questions:
\begin{itemize}
	\item How do recurrent, convolutional, and attention-based
	      architectures compare under identical conditions?
	\item What is the accuracy--latency trade-off for each family?
	\item Are these models robust to class imbalance, which is endemic
	      in real network datasets?
\end{itemize}

\subsection{Contributions}

This paper makes the following contributions:
\begin{enumerate}
	\item A rigorous head-to-head comparison of LSTM, CNN, and
	      Transformer models on two widely-used benchmark datasets.
	\item A lightweight Transformer variant (\textit{NetFormer})
	      designed for edge deployment with $<$500k parameters.
	\item An open-source reproduction package including pre-processing
	      scripts, model definitions, and training notebooks.
\end{enumerate}

The remainder of this paper is organised as follows.
Section~\ref{sec:related} reviews related work.
Section~\ref{sec:method} describes our methodology.
Section~\ref{sec:experiments} presents experimental results.
Section~\ref{sec:discussion} discusses findings and limitations.
Section~\ref{sec:conclusion} concludes the paper.

% ─────────────────────────────────────────────────────────────────────────────
\section{Related Work}
\label{sec:related}
% ─────────────────────────────────────────────────────────────────────────────

\subsection{Traditional Approaches}

Early IDS such as Snort~\cite{ref:roesch1999} rely on pattern matching
against known attack signatures. While highly precise for known
threats, they exhibit zero generalisation to novel exploits. Statistical
approaches, including $k$-Nearest Neighbour and Support Vector Machines,
improved generalisation but struggle with high-dimensional, high-rate
packet streams~\cite{ref:tsai2009}.

\subsection{Deep Learning for Network Security}

Recurrent architectures, particularly LSTMs, were among the first deep
models applied to IDS due to their natural handling of sequential data.
Kim \textit{et al.}~\cite{ref:kim2016} demonstrated that a two-layer
LSTM achieves 97.1\% accuracy on NSL-KDD, outperforming shallow
classifiers by a significant margin.

CNNs were later adapted for traffic classification by treating fixed-
length flow windows as one-dimensional signals~\cite{ref:wang2017}.
Their parallelism advantage over recurrent networks makes them
attractive for high-throughput inference.

Attention mechanisms and Transformer architectures, originally
developed for natural language processing~\cite{ref:vaswani2017},
have recently been applied to time-series anomaly
detection~\cite{ref:xu2021}. However, their applicability to
low-latency network monitoring remains under-explored.

% ─────────────────────────────────────────────────────────────────────────────
\section{Methodology}
\label{sec:method}
% ─────────────────────────────────────────────────────────────────────────────

\subsection{Datasets}

We evaluate on two standard benchmarks:

\textbf{NSL-KDD}~\cite{ref:tavallaee2009}: A refined version of the
KDD Cup 1999 dataset containing 125,973 training and 22,544 test
records across four attack categories: DoS, Probe, R2L, and U2R.

\textbf{CICIDS-2017}~\cite{ref:sharafaldin2018}: A modern dataset
generated in a controlled environment containing 2.8 million flow
records with 14 attack types including brute-force, XSS, and
infiltration attacks. Table~\ref{tab:datasets} summarises key
statistics.

\begin{table}[h]
	\caption{Dataset Statistics}
	\label{tab:datasets}
	\centering
	\begin{tabular}{@{} l r r r @{}}
		\toprule
		\textbf{Dataset} & \textbf{Records}   & \textbf{Features}
		                 & \textbf{Attack \%}                              \\
		\midrule
		NSL-KDD (train)  & 125,973            & 41                & 46.5\% \\
		NSL-KDD (test)   & 22,544             & 41                & 46.0\% \\
		CICIDS-2017      & 2,830,743          & 78                & 16.8\% \\
		\bottomrule
	\end{tabular}
\end{table}

\subsection{Pre-processing}

All categorical features are one-hot encoded. Continuous features
are normalised to zero mean and unit variance using statistics
computed on the training split only to prevent data leakage.
Missing values, present in 0.3\% of CICIDS-2017 records, are imputed
with the feature median.

Class imbalance is addressed using Synthetic Minority Over-sampling
Technique (SMOTE)~\cite{ref:chawla2002} on the training set, yielding
a balanced 50/50 split of benign and attack traffic.

\subsection{Model Architectures}

\subsubsection{LSTM Baseline}

Our LSTM model consists of two stacked LSTM layers with hidden
dimension $d = 128$, followed by dropout ($p = 0.3$) and a
fully-connected classification head:

\begin{equation}
	h_t = \text{LSTM}(x_t, h_{t-1}, c_{t-1})
\end{equation}

\begin{equation}
	\hat{y} = \sigma\!\left(W_o \cdot h_T + b_o\right)
\end{equation}

where $h_T$ is the final hidden state, $W_o \in \mathbb{R}^{2 \times d}$
and $b_o \in \mathbb{R}^{2}$ are learnable parameters.

\subsubsection{CNN Baseline}

The CNN model applies three parallel convolutional filters of widths
$\{3, 5, 7\}$ to each input window, followed by max-over-time pooling
and concatenation:

\begin{equation}
	z_k = \max_{t}\left(\text{ReLU}(W_k * x + b_k)\right), \quad
	k \in \{1,2,3\}
\end{equation}

\begin{equation}
	\hat{y} = \text{softmax}\!\left(W_c \cdot [z_1; z_2; z_3] + b_c\right)
\end{equation}

\subsubsection{NetFormer (Proposed)}

Our proposed \textit{NetFormer} replaces the recurrence and convolution
with multi-head self-attention:

\begin{equation}
	\text{Attention}(Q,K,V) =
	\text{softmax}\!\left(\frac{QK^\top}{\sqrt{d_k}}\right)V
\end{equation}

\begin{equation}
	\text{MultiHead}(Q,K,V) =
	\text{Concat}(\text{head}_1,\ldots,\text{head}_h)\,W^O
\end{equation}

We use $h = 4$ attention heads, model dimension $d_{\text{model}} = 64$,
two encoder layers, and a feed-forward inner dimension of 128, yielding
488k total parameters.

\subsection{Training Protocol}

All models are implemented in PyTorch and trained for 50 epochs using
the Adam optimiser with learning rate $\eta = 10^{-3}$ and weight decay
$\lambda = 10^{-4}$. A cosine annealing schedule reduces $\eta$ to
$10^{-5}$ by the final epoch. Batch size is 512. Early stopping with
patience 10 is applied on the validation loss.

Algorithm~\ref{alg:training} summarises the training procedure.

\begin{algorithm}
	\caption{Model Training Procedure}
	\label{alg:training}
	\begin{algorithmic}[1]
		\REQUIRE Dataset $\mathcal{D}$, model $f_\theta$, epochs $E$
		\STATE Split $\mathcal{D}$ into train / val / test (70/15/15)
		\STATE Apply SMOTE to training split
		\STATE Normalise features using training statistics
		\FOR{epoch $e = 1$ to $E$}
		\FOR{each mini-batch $(X, y) \in \mathcal{D}_{\text{train}}$}
		\STATE $\hat{y} \leftarrow f_\theta(X)$
		\STATE $\mathcal{L} \leftarrow \text{CrossEntropy}(\hat{y}, y)$
		\STATE $\theta \leftarrow \theta - \eta \nabla_\theta \mathcal{L}$
		\ENDFOR
		\STATE Evaluate on $\mathcal{D}_{\text{val}}$; update LR scheduler
		\IF{early stopping criterion met}
		\STATE \textbf{break}
		\ENDIF
		\ENDFOR
		\RETURN best $\theta$ by validation loss
	\end{algorithmic}
\end{algorithm}

% ─────────────────────────────────────────────────────────────────────────────
\section{Experiments}
\label{sec:experiments}
% ─────────────────────────────────────────────────────────────────────────────

\subsection{Evaluation Metrics}

We report Accuracy (ACC), Precision (PRE), Recall (REC),
$F_1$-score, and False Positive Rate (FPR):

\begin{align}
	\text{ACC} & = \frac{TP + TN}{TP + TN + FP + FN}          \\[4pt]
	F_1        & = \frac{2 \cdot \text{PRE} \cdot \text{REC}}
	{\text{PRE} + \text{REC}}                                 \\[4pt]
	\text{FPR} & = \frac{FP}{FP + TN}
\end{align}

\subsection{Results on NSL-KDD}

Table~\ref{tab:nslkdd} presents binary classification results on
NSL-KDD. \textit{NetFormer} achieves the best accuracy and $F_1$
while maintaining the lowest false positive rate.

\begin{table}[h]
	\caption{Binary Classification Results on NSL-KDD}
	\label{tab:nslkdd}
	\centering
	\begin{tabular}{@{} l c c c c c @{}}
		\toprule
		\textbf{Model}            & \textbf{ACC}  & \textbf{PRE}  & \textbf{REC}
		                          & \textbf{F1}   & \textbf{FPR}                              \\
		\midrule
		SVM~\cite{ref:tsai2009}   & 92.4          & 91.8          & 93.0         & 92.4 & 5.8 \\
		LSTM~\cite{ref:kim2016}   & 97.1          & 96.8          & 97.4         & 97.1 & 2.4 \\
		CNN                       & 97.5          & 97.2          & 97.8         & 97.5 & 2.1 \\
		\textbf{NetFormer (ours)} & \textbf{98.7} & \textbf{98.5}
		                          & \textbf{98.9} & \textbf{98.7}
		                          & \textbf{0.8}                                              \\
		\bottomrule
	\end{tabular}
\end{table}

\subsection{Results on CICIDS-2017}

Table~\ref{tab:cicids} shows multi-class results across all 14 attack
categories on CICIDS-2017. NetFormer again leads on macro-averaged
$F_1$, with particular strength on the low-frequency U2R and R2L
categories that challenge recurrent baselines.

\begin{table}[h]
	\caption{Multi-class Results on CICIDS-2017 (Macro Avg.)}
	\label{tab:cicids}
	\centering
	\begin{tabular}{@{} l c c c @{}}
		\toprule
		\textbf{Model}     & \textbf{ACC}  & \textbf{Macro-F1} & \textbf{FPR} \\
		\midrule
		LSTM               & 95.3          & 91.2              & 3.6          \\
		CNN                & 96.1          & 92.7              & 3.1          \\
		\textbf{NetFormer} & \textbf{97.4} & \textbf{94.8}     & \textbf{1.9} \\
		\bottomrule
	\end{tabular}
\end{table}

\subsection{Inference Latency}

Table~\ref{tab:latency} reports mean per-sample inference time
measured on an NVIDIA RTX 3060 GPU and an ARM Cortex-A72 CPU
(simulating edge deployment). NetFormer is $2.1\times$ faster than
the LSTM at batch size 1 due to the parallelisable attention
computation.

\begin{table}[h]
	\caption{Inference Latency ($\mu$s / sample)}
	\label{tab:latency}
	\centering
	\begin{tabular}{@{} l r r r @{}}
		\toprule
		\textbf{Model}     & \textbf{Params} & \textbf{GPU} & \textbf{CPU} \\
		\midrule
		LSTM               & 1.2M            & 42           & 310          \\
		CNN                & 0.9M            & 28           & 195          \\
		\textbf{NetFormer} & \textbf{0.5M}   & \textbf{20}  & \textbf{148} \\
		\bottomrule
	\end{tabular}
\end{table}

% ─────────────────────────────────────────────────────────────────────────────
\section{Discussion}
\label{sec:discussion}
% ─────────────────────────────────────────────────────────────────────────────

\subsection{Why Transformers Outperform RNNs Here}

Sequential models like LSTM process tokens one at a time, causing
gradient attenuation over long flow windows. Self-attention computes
pairwise interactions across the entire window in $O(n^2 d)$ operations,
allowing it to capture long-range dependencies between early and late
packets that are diagnostically significant for slow-scan attacks.

\subsection{Limitations}

Several limitations should be noted:
\begin{itemize}
	\item Both benchmark datasets are synthetic or semi-synthetic;
	      generalisation to live production traffic is unverified.
	\item We do not evaluate adversarial robustness --- a motivated
	      attacker may craft flows that evade the learned detector.
	\item SMOTE may introduce artefacts for highly imbalanced minority
	      classes such as U2R.
\end{itemize}

\subsection{Future Work}

Promising directions include: federated learning across multiple
network vantage points to improve generalisation without centralising
sensitive traffic; contrastive pre-training on unlabelled flows; and
formal verification of detection boundaries.

% ─────────────────────────────────────────────────────────────────────────────
\section{Conclusion}
\label{sec:conclusion}
% ─────────────────────────────────────────────────────────────────────────────

We presented a systematic comparison of LSTM, CNN, and Transformer
architectures for network anomaly detection. Our proposed
\textit{NetFormer} model achieves 98.7\% accuracy on NSL-KDD and
97.4\% on CICIDS-2017, with a false positive rate of 0.8\% and
1.9\% respectively --- improving over prior state-of-the-art whilst
reducing parameter count by $2.4\times$. The model's low latency
makes it viable for edge deployment, and we release all code and
pre-trained weights to support reproducibility.

% ─────────────────────────────────────────────────────────────────────────────
\begin{thebibliography}{99}
	% ─────────────────────────────────────────────────────────────────────────────

	\bibitem{ref:cisco2023}
	Cisco Systems, ``Cisco Annual Internet Report 2022--2027,''
	\textit{White Paper}, 2023.

	\bibitem{ref:liao2013}
	H.-J. Liao, C.-H. R. Lin, Y.-C. Lin, and K.-Y. Tung,
	``Intrusion detection system: A comprehensive review,''
	\textit{J. Network and Computer Applications}, vol. 36, no. 1,
	pp. 16--24, 2013.

	\bibitem{ref:ring2019}
	M. Ring, S. Wunderlich, D. Scheuring, D. Landes, and A. Hotho,
	``A survey of network-based intrusion detection data sets,''
	\textit{Computers \& Security}, vol. 86, pp. 147--166, 2019.

	\bibitem{ref:roesch1999}
	M. Roesch, ``Snort --- Lightweight intrusion detection for networks,''
	in \textit{Proc. USENIX LISA}, 1999, pp. 229--238.

	\bibitem{ref:tsai2009}
	C.-F. Tsai, Y.-F. Hsu, C.-Y. Lin, and W.-Y. Lin,
	``Intrusion detection by machine learning: A review,''
	\textit{Expert Systems with Applications}, vol. 36, no. 10,
	pp. 11994--12000, 2009.

	\bibitem{ref:kim2016}
	J. Kim, J. Kim, H. L. T. Thu, and H. Kim,
	``Long short term memory recurrent neural network classifier for
	intrusion detection,''
	in \textit{Proc. IEEE ICPADS}, 2016, pp. 1--4.

	\bibitem{ref:wang2017}
	W. Wang \textit{et al.},
	``HAST-IDS: Learning hierarchical spatial-temporal features using
	deep neural networks to improve intrusion detection,''
	\textit{IEEE Access}, vol. 6, pp. 1792--1806, 2017.

	\bibitem{ref:vaswani2017}
	A. Vaswani \textit{et al.},
	``Attention is all you need,''
	in \textit{Advances in Neural Information Processing Systems},
	vol. 30, 2017.

	\bibitem{ref:xu2021}
	J. Xu \textit{et al.},
	``Anomaly transformer: Time series anomaly detection with
	association discrepancy,''
	in \textit{Proc. ICLR}, 2022.

	\bibitem{ref:tavallaee2009}
	M. Tavallaee, E. Bagheri, W. Lu, and A. A. Ghorbani,
	``A detailed analysis of the KDD CUP 99 data set,''
	in \textit{Proc. IEEE CISDA}, 2009, pp. 1--6.

	\bibitem{ref:sharafaldin2018}
	I. Sharafaldin, A. H. Lashkari, and A. A. Ghorbani,
	``Toward generating a new intrusion detection dataset and intrusion
	traffic characterization,''
	in \textit{Proc. ICISSP}, 2018, pp. 108--116.

	\bibitem{ref:chawla2002}
	N. V. Chawla, K. W. Bowyer, L. O. Hall, and W. P. Kegelmeyer,
	``SMOTE: Synthetic minority over-sampling technique,''
	\textit{J. Artificial Intelligence Research}, vol. 16,
	pp. 321--357, 2002.

\end{thebibliography}

% ══════════════════════════════════════════════════════════════════════════════
\end{document}
% ══════════════════════════════════════════════════════════════════════════════
